% Methode: 1-1 1/2 Seiten: Datensatz, Modelle, Parameter (Folie 16).
\section{Methode}
\label{sec:method}

\subsection{Datensatz und Datenqualität}
Wir verwenden die 25 Hunderassen von Oxford-IIIT Pet~\cite{parkhi2012pets} (CC BY-SA 4.0) und lassen
die 12 Katzenrassen weg. Eine einzige Tierart macht die Aufgabe feinkörnig; haben wir Hunde gewählt, weil es
etwa doppelt so viele Bilder wie von Katzen gibt. Mit 25 Klassen liegen wir über den vom Modul empfohlenen 5--10
Klassen, jede Rasse hat aber noch etwa 200 Bilder. Der offizielle \texttt{trainval}-Split wird in
Training und Validation geteilt (10\,\% pro Klasse, stratifiziert, fester Split-Seed); der offizielle
\texttt{test}-Split wird einmal nach der Modellauswahl verwendet (\cref{tab:splits}).

\begin{table}[t]
  \centering
  \caption{Aufteilung der 25 Hunderassen in Training, Validation und Test.}
  \label{tab:splits}
  \small
  \begin{tabular}{lrl}
    \toprule
    Split & Bilder & Quelle \\
    \midrule
    Training   & 2\,242 & \texttt{trainval} ohne Validation \\
    Validation & 250    & 10\,\% pro Rasse, stratifiziert \\
    Test       & 2\,486 & offizieller \texttt{test}-Split \\
    \bottomrule
  \end{tabular}
\end{table}

\paragraph{Datenqualität.}
Wir haben alle 4\,978 Bilder geprüft. Die Klassen sind nahezu balanciert
(189--200 Bilder pro Rasse). Die Bildgröße liegt im Median bei $500\times375$\,px; 133 Bilder haben eine kürzere Seite unter
224\,px und werden hochskaliert. Der Hund bedeckt im Mittel 40\,\% des Bildes, in 123 Bildern weniger als
10\,\%. Exakte Duplikate gibt es nicht, daher wurde kein Bild entfernt. \Cref{fig:examples} zeigt ein
Beispiel pro Rasse.

\paragraph{Augmentierung.}
Unser Standard (\emph{light}) nutzt zufällige Ausschnitte (60--100\,\% des Bildes) und horizontales
Spiegeln; vertikales Spiegeln und Farbänderungen vermeiden wir, weil Orientierung und Fellfarbe
Rassemerkmale sind. \emph{strong} ergänzt RandAugment~\cite{cubuk2020randaugment} (zufällig gewählte Bildoperationen) und Random
Erasing~\cite{zhong2020erasing} (zufällig maskierte Bildbereiche), \emph{none} nutzt nur Skalierung und zentralen Zuschnitt.

\begin{figure}[t]
  \centering
  \resultfigure{paper/beispiele.png}{\linewidth}{Beispiel pro Rasse mit Maskenumriss}
  \caption{Je ein Trainingsbild der 25 Hunderassen (das erste der Rasse im Trainingssplit). Gelb: Umriss der
  Segmentierungsmaske des Tieres. Manche Rassen sehen sich sehr ähnlich, etwa American Pit Bull Terrier und
  Staffordshire Bull Terrier.}
  \label{fig:examples}
\end{figure}

\subsection{Modelle und Training}
Alle Modelle sind auf ImageNet~\cite{deng2009imagenet,russakovsky2015imagenet} vortrainiert und werden mit Hugging Face
Transformers~\cite{wolf2020transformers} auf PyTorch~\cite{paszke2019pytorch} geladen: ViT-B/16
(\texttt{google/vit-base-patch16-224}, auf ImageNet-21k vortrainiert und auf ImageNet-1k feinjustiert; 86\,M Parameter mit unserem Kopf für 25 Klassen), ResNet-50
(\texttt{microsoft/resnet-50}, 24\,M) und ConvNeXt-T (\texttt{facebook/convnext-tiny-224}, 28\,M). Wir
trainieren mit AdamW~\cite{loshchilov2019adamw} (Weight Decay 0,05) und Cosine-Schedule~\cite{loshchilov2017sgdr}
mit 10\,\% Warm-up, bis zu 10 Epochen mit
Early Stopping, Batchgröße 32 und $224\times224$ Pixeln Eingabe. Jede Konfiguration läuft mit drei Seeds;
wir berichten Mittelwert $\pm$ Standardabweichung auf dem Testsplit. Der Checkpoint wird auf dem
Validierungssplit ausgewählt.

\paragraph{Auswertung.}
Wir berichten Accuracy und Macro-F1. Einen Unterschied zwischen zwei Methoden interpretieren wir nur, wenn
die Differenz der Mittelwerte größer ist als das Doppelte der größeren Seed-Standardabweichung beider
Methoden (\emph{Schwelle}); alle Aussagen wie \glqq innerhalb der Streuung\grqq{} beziehen sich darauf. Mit drei
Seeds ist das ein grobes Kriterium und kein statistischer Test; die Schwelle ist konservativer als ein
gepaarter Vergleich je Seed (siehe Limitationen).

\subsection{Anpassungsmethoden}
\label{sec:methods}
\paragraph{Linear Probing.} Der ViT-Backbone ist eingefroren; trainiert wird nur der neue lineare
Klassifikator auf dem \texttt{[CLS]}-Token (Lernrate $10^{-3}$).

\paragraph{LoRA.} Für ein eingefrorenes vortrainiertes Gewicht $W\in\mathbb{R}^{d_\text{out}\times
d_\text{in}}$ lernt LoRA~\cite{hu2022lora} $A\in\mathbb{R}^{r\times d_\text{in}}$ und
$B\in\mathbb{R}^{d_\text{out}\times r}$ mit Rang $r\ll d$:
\begin{equation}
  h = W x + \frac{\alpha}{r}\, B A x .
  \label{eq:lora}
\end{equation}
In \cref{eq:lora} wird $B$ mit Nullen initialisiert, sodass das Training beim vortrainierten Modell
beginnt. Wir passen die Query- und Value-Projektionen aller 12 Blöcke mit $r\in\{4,8,16\}$ an und für
$r=8$ zusätzlich alle linearen Layer jedes Blocks (Key, Attention-Ausgabe, beide MLP-Layer). Wir setzen
$\alpha=2r$, verwenden Dropout 0,1 im LoRA-Pfad und trainieren zusätzlich den Klassifikationskopf (Lernrate
$5\cdot10^{-4}$). Nach dem Training wird $\frac{\alpha}{r}BA$ zu $W$ addiert; das angepasste Modell hat
damit dieselbe Architektur und dieselben FLOPs wie der ursprüngliche ViT.

\paragraph{Volles Fine-Tuning.} Alle Parameter werden trainiert (Lernrate $5\cdot10^{-5}$); das ist die
Referenz, mit der die effizienten Methoden verglichen werden. Auch die CNNs werden vollständig
angepasst. Gemäß Vorgabe des Dozenten werden sie nicht von Grund auf trainiert, sondern starten mit vortrainierten
ImageNet-Gewichten und werden nur an unsere Aufgabe angepasst (Transfer Learning); das Training ohne
Pretraining in \cref{sec:why} dient ausschließlich als Vergleich.

\paragraph{Lernraten.} Die Lernraten der ViT-Methoden wurden einmal pro Methode festgelegt und nicht
optimiert. Unsere ersten CNN-Läufe verwendeten die
ViT-Lernrate $5\cdot10^{-5}$ und waren deutlich untertrainiert (Validation-Loss 1,35 bei ResNet-50). Daher
haben wir die CNN-Lernrate auf dem \emph{Validierungssplit} aus $\{5\cdot10^{-5}, 2\cdot10^{-4},
5\cdot10^{-4}, 10^{-3}\}$ mit je drei Seeds gewählt; $2\cdot10^{-4}$ war für beide CNNs am besten
(\cref{tab:lr}).

\begin{table}[h]
  \centering
  \caption{Wahl der CNN-Lernrate auf dem Validierungssplit (Mittelwert $\pm$ Std über 3 Seeds; Loss:
  Cross-Entropy; fett: gewählte Lernrate).}
  \label{tab:lr}
  \small
  \resizebox{\linewidth}{!}{\begin{tabular}{lcccc}
    \toprule
    & \multicolumn{2}{c}{ResNet-50} & \multicolumn{2}{c}{ConvNeXt-T} \\
    Lernrate & Val.-Acc.\ (\%) & Val.-Loss & Val.-Acc.\ (\%) & Val.-Loss \\
    \midrule
    $5\cdot10^{-5}$ & $91{,}47\pm1{,}29$ & 1,349 & $96{,}40\pm0{,}69$ & 0,303 \\
    $2\cdot10^{-4}$ & $\mathbf{97{,}47\pm0{,}46}$ & 0,146 & $\mathbf{97{,}47\pm0{,}46}$ & 0,112 \\
    $5\cdot10^{-4}$ & $96{,}40\pm1{,}06$ & 0,110 & $96{,}13\pm0{,}83$ & 0,157 \\
    $10^{-3}$ & $96{,}13\pm1{,}80$ & 0,154 & $94{,}80\pm0{,}80$ & 0,190 \\
    \bottomrule
  \end{tabular}}
\end{table}
